\section{序論：反標準heartと二種類の根現象}\label{sec:intro}
滑らかな可換del Pezzo曲面では，点のblow-up，反標準環，例外列から作るhelixは
同じcoherent sheaf圏を回復する．非可換の場合には，同じ導来圏に異なる標準heartが現れ，
導来同値の存在だけでは構造対象や反標準切断の同定は決まらない．
本稿の問いは，\emph{指定された構造対象と反Serre作用がどのheartを選ぶか}である．
その答えには，非零指数の共鳴に伴うheartのtiltと，零指数壁に伴う対象の収縮という
二つの異なる現象が現れる．

構造対象を$e$，Serre関手を$\Sfun$，反Serre作用を$s=\Sfun^{-1}[2]$とし，
$R(e,s)=\bigoplus_{n\ge0}\Hom(e,s^ne)$を反標準orbit環と呼ぶ．
名称は構成の由来を表し，\emph{VdB heart}はVan den Berghのraw blow-up，
\emph{Rains heart}はgluingとblowdown表示で指定されるt構造，
\emph{AS heart}はArtin--Schelter正則$\mathbb Z$-代数，
\emph{反標準heart}は指定された$R(e,s)$の$qgr$から得る．
四名称が四つの相異なるabel圏の存在を意味するわけではない．
圏の構成，導来同値，そのt-exactness，構造対象・境界・標識の保存を分けて扱う．

\subsection{genuine領域：heartの選択と有限次元表現}
滑らかな射影有理曲面$S$の$H=-K_S$がampleである場合をgenuineと呼ぶ．
滑らかな楕円境界$E\in|H|$と$q\in\Pic^0(E)$による標準Rains gluingを固定すると，
比較の基本図式は次のようになる．
\begin{equation}\label{eq:overview}
\begin{tikzcd}[column sep=large,row sep=large]
\AV \arrow[r,"\text{同じblow-up表示}","\simeq"'] & \AR
\arrow[d,phantom,"\text{HRS tilt}" description] \\
\AAS \arrow[r,"\text{指定AS実現}","\simeq"'] & \AAC
\end{tikzcd}
\end{equation}
右側の縦の表示はabel同値ではなく，指定導来同値の下でのheartの関係である．
下側にはfull very strong helixとrolled-up同定などの入力があり，
任意のpotentialに対して存在する矢とは解釈しない．

\begin{theorem}[genuine領域の比較と表現論]\label{thm:summary}
以下はそれぞれに記した仮定の下で成り立つ．
\begin{enumerate}
\item 滑らかなgenuine $S$，滑らかな楕円境界$E$，任意の$q\in\Pic^0(E)$からの
標準gluingでは，自然反標準導来切断$\Phi$はdg同値であり，Rains heart上の振幅は$[0,1]$である．
従って
\[
 \Phi^{-1}(\AAC)=\langle\cF[1],\cT\rangle,\qquad
 \cT=\ker H^1\Phi,\quad\cF=\ker H^0\Phi
\]
となる（定理\ref{thm:derived}，\ref{thm:globaltilt}）．
\item さらに$q$が非torsionなら，$\cF$は境界非交叉の有限長Serre圏$\cZ$に含まれる．
その実在単純根は有限ADE図式$\Gamma$をなし，
\[
 \Phi\text{ がt-exact}\iff(e,s)\text{ がAZ-ample}
 \iff\chi(e,S_i)>0\quad(\text{全ての実在単純 }S_i\in\cZ)
\]
である．負指数単純の有限回のtiltが標的heartを回復する
（定理\ref{thm:globaldefect}，\ref{thm:Eulerchamber}，\ref{thm:globalADEtilts}）．
\item $\cZ$と標的の有限次元heartは，それぞれ$\operatorname{mod}_{\rm fd}\Pi(\Gamma)$とabel同値である．
$A=R/(g-1)$，標的単純の次元$m_i=\chi(e,\beta_i)>0$に対して，
全有限次元$A$加群は一つの有限次元商
\[
 A/I\simeq\End_{\Pi(\Gamma)}\left(\bigoplus_iP_i^{\oplus m_i}\right)
\]
を経由する（定理\ref{thm:ADE-abelian}，系\ref{cor:ADE-affine-quotient}）．
これは162次元の頂点除去cornerとは別の代数である．
\item $D=A/I$について，非被約基底を含む古典的な表現スキームは
$\operatorname{Rep}_n(A)\simeq\operatorname{Rep}_n(D)$となる．
表現stackは$\sum_i m_id_i=n$を満たすpreprojective表現stackの和である
（定理\ref{thm:finite-rep-stack}）．$D$は自己入射的で，
Frobeniusとなる条件は$m_i=m_{\nu(i)}$である（命題\ref{prop:quotient-Frobenius}）．
\item 同じ初期heart，中心，gluing bimodule，blow-down関手を同定したVdB表示とRains表示では，
誘導される同値はt-exactである（定理\ref{thm:RV}）．同じ最後のraw blow-upについては
section-model関手を自然反標準関手へ同定できる
（定理\ref{thm:VdBcanonical}，系\ref{cor:fourfunctors}）．
\item 三次曲面のAS helixのfoldingと反標準orbitの同定，命題\ref{prop:corner}の有限性入力の下では，
自然corner restrictionがabel同値となる必要十分条件は$\dim B/BpB<\infty$である．
規格化decupleでは全27theta係数から入力の$L_1$六点を復元でき，頂点除去cornerは$L_0$に依存しない．
指定されたfull very strong AS実現とrolled-up同定の下で，AS--反標準比較そのものがt-exactである
（系\ref{cor:ASac}，定理\ref{thm:theta-potential}，\ref{thm:original-six-points}，\ref{thm:normalized-abelian}）．
\end{enumerate}
\end{theorem}
実在する負指数根sheafは，導来切断で消えるのでなく次数1へ移る．
それが反標準AZ amplenessの失敗を検出する（系\ref{cor:negative}）．
一般の対象の総Euler値だけではtorsion pairは判定できず，
全ての非零商または部分対象のEuler値を用いる（命題\ref{prop:numericaltorsion}）．

整数倍点$p_i=N_iq$の実在配置では，単純tiltは整数列の整列として記述される．
一つの負指数根だけを持つ配置と$E_6$の混合配置を第\ref{sec:single-root}節までに具体化する．
$E_6$の普遍有限商は例によって970次元となり，そのbasic代数の156次元とも区別される．
生成元ごとの作用行列を計算したことと，忘却関手が指定する商の同型型を決めたことは異なる．

\subsection{最大共鳴と零指数壁}
固定した$E,q$について，相異なる正整数倍点は次数1--6の全根を非零指数で共鳴させる．
任意の厳密な整数順序を課してもその配置はZariski稠密である
（定理\ref{thm:maximal-density-main}）．整列定理と合わせると，
指定heartが反標準heartと一致する点も一致しない点も稠密であり，
その一致条件は構成可能でない（系\ref{cor:nonconstructible-texact}）．
これは点ごとに指定したheartの結論で，代数的な相対t構造の存在や全曲面モジュライの成分分類ではない．

一方，$H=-K_S$がnefかつbigでampleでない弱del Pezzoでは，
反標準収縮$f:S\to\overline S$の核が存在する．
古典的な$(-2)$曲線$C$の$i_*\OO_C(-1)$は全反標準orbitから見えず，
零指数での対象の消失は，genuine領域のtiltと異なる現象である．

\begin{theorem}[弱gluingのcoherent収縮と商圏の記述]\label{thm:weak-summary}
$S$を滑らかな射影有理弱del Pezzo曲面，$E\in|-K_S|$を滑らかな楕円曲線とし，
任意の$q\in\Pic^0(E)$による標準gluing $\cD_q$を取る．
\begin{enumerate}
\item $R_q$は左右noetherianで，各$X\in\cD_q$の次数付きExt
$\bigoplus_{n\ge0}\Hom(e,s_q^nX[j])$は有限生成で，非零の$j$は有限個である．
従ってcoherentな自然dg関手$\Phi_q:\cD_q\to\Db(\qgr R_q)$が存在する．
\item その全核は
\[
 \ker\Phi_q=\mathcal N_q
 =\iota_q\{G\in\Perf(S):Rf_*G=0\}
\]
である．ここで$\mathcal N_q$は全反標準orbitに直交するdg部分圏で，
$\iota_q$は共通gluing成分$\OO_S^\perp$の包含である．
\item $(\cD_q/\mathcal N_q)^\natural$は，$\OO_{
\overline S}^{\perp}\subset\Db(\coh\overline S)$にexceptional対象を
$\RHom_E(L\bar i^*(-),q)$でglueした特異gluing $\overline{\cD}_q$と同値である．
$q=\OO_E$ではこれは$\Db(\coh\overline S)$である．
\end{enumerate}
証明は定理\ref{thm:weak-orbit-kernel}，\ref{thm:weak-coherent-sections}，
\ref{thm:singular-gluing-target}に与える．
\end{theorem}
この定理は，coherent有限性を追加仮定として課すことなく標準弱gluingへ適用できる．
ただし一般の$q$について$\overline{\cD}_q\simeq\Db(\qgr R_q)$までは結論しない．
その比較は自然関手$\overline\Phi_q$の完全忠実性と本質的全射性に帰着する．
非genuineの場合，$\overline{\cD}_q$には無限個の高次数に自己Extを持つ対象があり，
source $\cD_q$のpropernessを商に移してはいけない（系\ref{cor:singular-gluing-unbounded}）．
さらに古典的特異性圏$D_{\mathrm{sg}}(\overline S)$は共通部分圏$\overline{\cB}$の商として残り
（命題\ref{prop:duval-singularity-quotient}），Du Val点$x_i$のideal sheafの像は
全$n\ge3$で高次Ext次元$4\delta_{ij}$を持ち，
gluingパラメータに依存しない（定理\ref{thm:duval-high-ext}）．
さらに特異点idealとgluingのexceptional対象の間には
$\RHom(J_{x,q},\bar e)\simeq R\Gamma(E,q)$が成立し，
非自明$q$では両側直交する一方，$q=\OO_E$では次数$0,1$に各一次元のExtが現れる
（命題\ref{prop:duval-ideal-boundary-gluing}）．これは古典的局所自己Extの$q$不変性と
例外成分へのgluingの$q$依存性を区別するものである．
また，相異なるDu Val点のideal間の導来Homは$k[-1]$である
（命題\ref{prop:duval-offdiagonal-ext}）．
非自明$q$では，これらのidealと$\bar e$が生成する有限部分圏は直和に分かれるが，
$\overline{\cD}_q$全体の分解は主張しない
（系\ref{cor:duval-local-splitting}）．
重要なことに，$J_{x,q}$の反標準$qgr$への像は実際にcoherent対象であり，
その切断加群は$H_R(t)-H_M(t)=1/(1-t)$を満たす
（定理\ref{thm:duval-point-hilbert}）．これは$qgr$側で非perfect対象を
調べるための具体的な出発点である．
全比較が完全忠実なら非自明$q$で
$\RHom_{\qgr R_q}(\mathcal I_{x,q},\pi R_q)$の全次数の消滅が必要となる
（系\ref{cor:duval-target-hom-obstruction}）．
さらに一つの$A_1$壁では$uv-w^2$の明示的な一周期matrix factorizationから
局所特異性圏のidealの安定自己準同型代数を
$M_2(k[\eta,\eta^{-1}])$（$\deg\eta=1$）と同定する
（命題\ref{prop:A1-stable-end}）．
これは普通のExtの低次数の積でも，非自明$q$の反標準$qgr$の計算でもない．
一方，反標準orbitから生成される完全部分圏
$\mathcal T_q=\thick\{s_q^ne:n\in\mathbb Z\}$上では，
弱領域でも自然関手が$\thick\{\pi R_q(n)\}$への同値を与える
（定理\ref{thm:orbit-core-equivalence}）．
非genuineならこの完全部分圏は商$\mathcal Q_q$全体を生成せず，
残余商$\mathfrak S_q$は非零である．
$q=\OO_E$では$\mathfrak S_q$は古典的な$D_{\mathrm{sg}}(\overline S)$に一致する
（定理\ref{thm:residual-singularity}）．

\subsection{先行研究と寄与の範囲}
Van den Bergh\cite{VdB}は非可換blow-upと，一般位置の下でのsection algebraへの導来比較を構成した．
Rains\cite{Rains,Filtered}はそのgluingとt構造，effective roots，数値的なnef/ample coneを解析した．
数値的amplenessと指定pairのAZ amplenessは定義から同じではない．
本稿のVdB--Rains一致は既知のpseudo-canonicalな認識の適用であり，
追加的な問題は\emph{自然な反標準関手が選ぶheart}を実在対象で決定することである．
Rainsの2025年monograph\cite{RainsBook}では，旧稿のLemma 4.9，Proposition 4.17，
Theorem 4.30，Theorem 9.10，Corollary 10.16はそれぞれ
Lemma 11.10，Proposition 11.19，Theorem 11.31，Theorem 16.10，Corollary 17.17に対応する．
単純根sheafの存在にはProposition 18.27直前の記述を用い，nef仮定付きの同命題自体を根へ誤適用しない．

Abdelgadir--Okawa--Ueda\cite{AOU}の三次曲面型構成とOkawa--Ueda\cite{OU,OUHelices}の
helix理論はAS比較の入力である．本稿では実際の右加群shiftから境界辞書
$L=L_1$，$q=L_0L_1^{-1}$を運び，構造対象を保存する規格化mutationを固定する
（第\ref{sec:boundarydata}節，定理\ref{thm:AOU-pointed}）．
AOUの印刷座標の異なるshift規約を混在させず，未変更の全座標への翻訳を確定したとも主張しない．
指定$\Psi_+$のt-exactness，actual theta係数と入力六点の一致，$L_0$非依存を
この規格化の下で証明する．elliptic Cauchy恒等式自体は既知\cite{EllipticCauchy}である．
係数トーラスの十一条件と一致領域は定理\ref{thm:AOUexactlocus}，系\ref{cor:coefficientopen}に記す．

HRS理論\cite{CHZ}，Kleinian圏のWeyl反射\cite{BridgelandKleinian}，
length heartの単純tilt\cite{BCPW}，ADE zigzagの形式性\cite{LiuWangZigzag}と
拡大閉包の復元\cite{KellerAinf}，Dynkin torsion classの分類\cite{Mizuno}は既知の機構である．
本稿ではこれらを二つの実在heartへ適用し，Euler次元から元のアフィン代数の
\emph{特定の非basicな有限商}を取り出す．自己入射性・Nakayama置換\cite{AiharaMizuno}と
matrix Hilbert公式\cite{EtingofRains}も既知入力として区別する．
GDAHAとmultiplicative preprojectiveを通じたquantized del Pezzoの先行研究\cite{EOR}は，
境界や表示を確認せずに本稿の滑らかな楕円gluingへ同一視しない．

弱領域ではRainsの一般有理曲面のgluingを出発点とし，
境界の正次数性とKeelerのnoetherian性\cite{KeelerAmple}からcoherent導来切断を構成する．
古典的な反標準収縮\cite{DolgachevWalls}に対する
Bondal--Kapranov--Schechtmanのfiber次元高々1の局所化定理\cite[Theorem 2.14]{BKSWalls}を用い，
dg商\cite{DrinfeldWalls}を特異モデル上のgluingに同定する．
これらの一般定理自体を新しい結果とはせず，固定された$(S,E,q,e,s)$の
全orbit核・coherent関手・収縮圏を接続することを本稿の結論とする．
与えられた収縮から始まる非可換収縮環\cite{DonovanWemyssWalls}は，
任意の非可換曲面上の収縮射の存在を代替するものではない．

\subsection{仮定・未解決点と第2論文の役割}
自然導来\emph{同値}とHRS比較はgenuine性を使う．実在ADE根によるheart選択と
アフィン有限商には非torsion性を加える．AS比較には指定helixとrolled-up同定を加える．
これに対しcoherent反標準\emph{関手}と全核・特異gluing表示は弱del Pezzoまで成り立ち，
$q$の位数制限を必要としない．任意の非可換曲面，特異楕円境界，相対t構造への拡張は含まない．

第2論文\texttt{dp-hearts2.tex}は，符号別Ext・混合拡大の詳細，固定dg圏のframed形式変形，
三次曲面cornerの完全Hilbert関数と162次元平坦族・厳密計算を扱う．
良いgeometric helixのcorner有限性自体は既知一般論\cite{dTdVV}でも説明できる．
本稿は自然比較に必要な有限性と六点復元を自立して証明し，細かいHilbert計算を重複させない．
零指数壁での一般$q$の$qgr$同定，指定heartでの収縮対象の位置，局所完成代数，
曲面族の底変更・降下は第\ref{subsec:coherent-future}節の課題として明記する．
これらの未解決点と，本稿で証明するpointwiseな結論を混同しない．

\section{記法，realization，反標準環}\label{sec:setup}
基礎体は$k=\mathbb C$とし，圏は$k$-線形で適合するdg enhancementを用いる．
sourceの滑らかなgluing $\cD_q$は冪等完備でproper，すなわち各RHomの全cohomologyが有限次元である．
零指数壁の商$\mathcal Q_q$や特異gluing $\overline{\cD}_q$にはpropernessを仮定しない．
有界heart$\cH\subset\cD$からのrealization
$\operatorname{real}_{\cH}:\Db(\cH)\to\cD$が同値となる場合を\emph{derived type}と呼ぶ．
heartが有界であるだけではこの同値は従わない．

Serre関手を$\Sfun$，$\theta=\Sfun[-2]$，$s=\theta^{-1}=\Sfun^{-1}[2]$と書く．
$s$を単にAS代数の次数1 shiftと同一視しない．一般の$\mathbb Z$-代数には
一周期的なshiftの指定自体がなく，helixの一周期は複数頂点を通ることがある．
derived typeのheartを$s$が両方向で保つなら
\[
\Ext^i_{\cH}(M,N)=\Hom_{\cD}(N,s^{-1}M[2-i])^*=0\quad(i>2)
\]
なので，$\cH$のhomological dimensionは高々2である．
$\OO$がexceptionalなら$\Ext^2(\OO,s^{-1}\OO)=k$であり，ちょうど2となる．

\begin{definition}[構造対象つき比較]
$(\cH,\OO,s)$の間の\emph{pointed同値}とは，構造対象の同型と$s$との自然な両立同型を
指定した同値をいう．境界restrictionやmarkingの保存は，必要な箇所で別に指定する．
\end{definition}
構造対象を$e=\OO$と書き，右加群の規約で
\[
R=R(e,s)=\bigoplus_{n\ge0}\Hom(e,s^ne),
\qquad ab=s^{\deg b}(a)\circ b
\]
と置く．局所有限右noetherian性を確認した範囲で
$\AAC=\qgr R=\gr R/(\text{有限次元加群})$と定める．
$\AR$と$\AAC$を同じ三角圏のheartと呼ぶには，まず$\Phi$を構成してtransportする必要がある．

\begin{definition}[AZ ample pair]
noetherian abel圏のpair $(e,s)$がAZ-ampleであるとは，任意の対象が
有限個の$s^{-n}e$（$n>0$）の直和の商であり，任意の全射$M\twoheadrightarrow N$に対して
$\Hom(e,s^nM)\to\Hom(e,s^nN)$が$n\gg0$で全射となることをいう\cite{AZ}．
\end{definition}
高い全Extの消滅をこの二条件だけから無条件に推論しない．
またhelixの全頂点による生成性と，一つの頂点の$s$-orbitによる生成性は別である．

\section{比較の一般論：振幅1，ample orbit，corner}\label{sec:formal}
\subsection{振幅1の同値とHRS tilt}
\begin{proposition}\label{prop:hrs}
$F:\Db(\cA)\simeq\Db(\cB)$をenhancementと両立する同値，$G=F^{-1}$とし，
$F(\cA)\subset D^{[0,1]}(\cB)$とする．このとき
\[
\cT=\ker H^1F,\qquad\cF=\ker H^0F
\]
は$\cA$上のtorsion pairであり，
\[
G(\cB)=\{X:H^{-1}_{\cA}X\in\cF,\ H^0_{\cA}X\in\cT,
\ H^i_{\cA}X=0\ (i\ne-1,0)\}=\langle\cF[1],\cT\rangle.
\]
これはderived typeのheartである．$F$がt-exactであることと$\cF=0$は同値である．
\end{proposition}
\begin{proof}
有界t構造のaisleを有限拡大とshiftで記述すると，振幅仮定と直交性から
$G(\cB)\subset D^{[-1,0]}(\cA)$を得る．
$M\in\cA$，$B_i=H^i(FM)$に対しtruncation triangleへ$G$を施すと
$GB_0\to M\to GB_1[-1]$となる．長完全列から
$H^{-1}_{\cA}(GB_0)=H^0_{\cA}(GB_1)=0$．よって
\[
0\longrightarrow GB_0\longrightarrow M\longrightarrow
H^{-1}_{\cA}(GB_1)\longrightarrow0
\]
は$\cT$と$\cF$への分解である．$T\in\cT,U\in\cF$では
$FT\in\cB$，$FU\in\cB[-1]$なので$\Hom(T,U)=0$．
逆に二つのcohomologyが表示条件を満たす$X$は$FX\in\cB$となり，heartの式を得る．
このheartのrealizationは$G$を通して標準の$\Db(\cB)$と同定されるため同値である．
$\cF=0$ならtorsion分解から$\cT=\cA$で，両heartは一致する．逆は明らか．
\end{proof}
\begin{remark}
$\cF\ne0$の対象は，$F$で消えるのでなく$H^1(FM)[-1]\ne0$になる．
抽象的なtorsion pairからHRS heartを作るだけではderived type性は従わない\cite{CHZ}．
\end{remark}

\subsection{反標準orbitによるheartの一意性}
\begin{proposition}\label{prop:unique}
$\cH\subset\cD$はderived typeの有界heart，$s$は$\cH$を両方向で保ち，
$(e,s)$はAZ-ampleとする．さらに各$M\in\cH$について，すべての$p>0$に共通の閾値で
$\Ext^p(e,s^nM)=0$ $(n\gg0)$を仮定する．このとき
\[
\cH=\{X\in\cD:\Hom(e,s^nX[i])=0\quad(i\ne0,\ n\gg0)\}.
\]
従ってこの条件を満たすheartは固定した$(\cD,e,s)$に対し高々一つである．
\end{proposition}
\begin{proof}
有限個の$H^j_{\cH}X$に高次消滅を同時に適用し，hyper-Ext spectral sequenceを退化させると
$\Hom(e,s^nX[i])=\Hom(e,s^nH^i_{\cH}X)$となる．
非零$M$について右辺がすべての十分大きい$n$で零となることは，任意に大きい$N$で
$s^NM$を負twistの有限直和から生成するAZ条件に反する．従って指定した消滅は
$H^i_{\cH}X=0$ $(i\ne0)$と同値である．
\end{proof}
\begin{corollary}\label{cor:ampleexact}
二つのheartが命題\ref{prop:unique}の条件を満たすなら，pointed導来同値はt-exactである．
同じ条件のheartの非零$s$-固定対象$M$は$\chi(e,M)>0$を満たす．
$s^rM\simeq M$なら少なくとも$\sum_{a=0}^{r-1}\chi(e,s^aM)>0$である．
\end{corollary}
\begin{proof}
第一の主張はheartの認識式による．固定対象では高次Extが消え，
$\chi(e,M)=\dim\Hom(e,M)$となる．この値が零なら全orbitからのHomが零となってAZ生成性に反する．
周期対象には$\bigoplus_{a=0}^{r-1}s^aM$を用いる．
\end{proof}

\subsection{idempotent cornerの判定}
\begin{proposition}\label{prop:corner}
$B$は局所有限右noetherian非負次数付き代数，$p\in B_0$はidempotentとする．
$R=pBp$が右noetherian，$Bp$が右$R$有限生成，$pB$が左$R$有限生成であると仮定する．
自然な$(-)p:\qgr B\to\qgr R$が同値である必要十分条件は$\dim_k B/BpB<\infty$である．
\end{proposition}
\begin{proof}
$Bp$の有限生成性から$Mp$は有限生成であり，$(-)p$はexactかつ有限次元加群を保つ．
左随伴$L(N)=N\otimes_RpB$は有限生成加群を有限生成加群に送り，
$pB$の左有限生成性から有限次元加群も有限次元に送る．
$B/BpB$が有限次元なら，counit $Mp\otimes_RpB\to M$の核・余核は$p$で消される．
それらは右noetherian性により有限生成の$B/BpB$加群なので有限次元である．
$L$を短完全列へ施した際の左端の核も同じ議論で有限次元となり，$L$は商圏上exactとなる．
$(N\otimes_RpB)p=N$とcounitが逆同値を与える．
逆に$(-)p$が同値なら$\pi(B/BpB)$は零に送られるので零であり，有限生成性から有限次元である．
\end{proof}
